\section{Related Work}

\paragraph{Generative agents and behavior simulation.}
Generative agent systems use language models, memory, reflection, and planning to simulate believable behavior in interactive environments \cite{park2023generative}, and broader surveys organize this design space around perception, memory, planning, and action \cite{wang2024agentssurvey}. Our setting differs: rather than populating a sandbox with synthetic agents, we replay one real user's historical GPT interactions and evaluate next-action prediction on held-out contexts.

\paragraph{Personalized language models.}
Personalization benchmarks such as LaMP evaluate whether language models can adapt to user-specific histories \cite{salemi2023lamp}, and recent surveys distinguish personalization by granularity, technique, and evaluation protocol \cite{zhang2024personalizationsurvey}. Most personalization work asks whether a model produces a better answer for a personalized task; our replay task asks whether an explicit user representation predicts the behavioral sequence of one user's conversation.

\paragraph{Memory-augmented assistants and personal digital twins.}
Long-term memory systems for LLMs persist facts, preferences, and episodes across sessions. MemoryBank studies memory for companion-like interaction \cite{zhong2023memorybank}, MemGPT frames memory management as a virtual-context problem \cite{packer2023memgpt}, and surveys emphasize that memory evaluation remains an open problem \cite{zhang2024memorysurvey}. Digital-twin discussions share the ambition of persistent personal models but often mix memory, goals, and simulation into a broad identity claim. We align with this motivation but treat ``digital twin'' as a measurable replay problem, explicitly separating persona, dynamic state, and memory to test which component improves held-out behavior prediction.

\paragraph{User simulation and social interaction.}
User-simulation work produces realistic user behavior in dialogue, task, or social settings. SOTOPIA evaluates language agents in open-ended social interactions \cite{zhou2023sotopia}, with follow-up work training agents through interactive social learning \cite{wang2024sotopiapi}. Language models have also been used to simulate human samples in social-science settings \cite{argyle2023outofone}. Such simulators are valuable for evaluation and agent training, but they do not show that a model can replay a specific person's historical trajectory. Our results support a related concern: static persona facts alone are not the dominant source of held-out replay gain, and a stronger semantic context comparator can exceed a learned state router.

\paragraph{Positioning.}
The closest existing directions provide believable agents, personalized task models, persistent memory, or simulated populations. This paper targets a different evaluation unit: one user's historical conversation trajectory with proxy next-action labels. The contribution is not a general-purpose agent architecture, memory store, or population simulator. It is a state-aware replay formulation for personal chat logs with a controlled comparison showing that apparent state-routing gains can disappear under stronger lexical comparators. This deliberately conservative positioning is useful when a researcher wants an auditable proxy for personal behavioral continuity before making stronger claims about digital twins.
